% ============================================================
%  MODULO 20 — Podcast por modulo: indice de episodios
%  SPEC R686 — um audio curto pt-BR por modulo (fila progressiva)
% ============================================================
\chapter{Podcast Por M\'odulo: Ou\c{c}a Cada Parte}

\roteirodidatico
{Se um capitulo ficou denso, ou\c{c}a o episodio dele. Cada modulo do livro tem um audio curto em portugues com detalhe, funcionalidade, justificativa, intera\c{c}\~oes, agentes, specs e hooks.}
{\emph{Episodio por modulo} aprofunda uma parte; \emph{consolidado} apresenta o todo. Ou\c{c}a o todo antes das partes.}
{Escolha o modulo travado, ou\c{c}a o episodio com a ficha aberta e pause para conferir arquivo e comando.}
{Audio resume fontes do caderno; prova esta no texto e no DOI. Fila de gera\c{c}\~ao respeita a cota do plano.}
{Que modulo voce ouve hoje, e que detalhe dele voce reconta depois?}

\casoguiado{Bruno destrava o roteamento}
{Bruno trava no roteamento por aten\c{c}\~ao. Ouve os seis minutos do modulo 02 com o fluxograma M6 aberto, pausa no Jaccard, refaz a conta 2 sobre 4 e volta ao codigo. Destrava sem reler o capitulo inteiro.}

\section{Episodios prontos e fila}

\elemento{\'Indice de episodios por modulo}{\metainfo{ID}{PODM-01} \hfill \metainfo{Idioma}{pt-BR curto} \hfill \metainfo{Pasta}{output/podcast/modulos/}}

\begin{descricao}
Prontos: 02 orquestra\c{c}\~ao E1-E7, 03 memoria MCI, 04 qualidade SDD e TDD, 05 raciocinio, 07 integra\c{c}\~oes com MCP e executores, 10 catalogo com 238 agentes em 14 familias, 13 ativa\c{c}\~ao com trilha auditavel. Gerando: 01 funda\c{c}\~oes e 06 ci\^encia. Fila apos renovar a cota: 00a prologo, 00b jornada, 01b historia, 01c instala\c{c}\~ao, 01d arquitetura, 08 ecossistema, 09 infra, 09a comparada, 12 fichas nucleo, 14 lotes, 15 epilogo, 16 atlas, 17 E1-E7 leigo, 18 nivel 0 e 19 podcast. Cada episodio traz detalhe, funcionalidade, justificativa com DOI, intera\c{c}\~oes, agentes e subagentes, specs, hooks e ferramentas do modulo. Para gerar a fila, crie o audio no caderno por modulo com idioma pt-BR e foco do modulo, depois baixe para a pasta de modulos.
\end{descricao}

\begin{funcao}
Transformar cada capitulo denso em conversa curta, sem perder rastro: o episodio aponta a ficha, a ficha aponta o arquivo.
\end{funcao}

\begin{aviso}
Limite dos episodios. Audio gerado pode simplificar nomes e numeros; o numero que vale e o recalculado no checkout. Cota do plano limita a fila: episodios restantes saem apos renovar a janela.
\end{aviso}
